\section{Discussion}\label{discussion}

\textcolor{red}{\subsection{Reframing the Drake Equation's Longevity Term}\label{reframing-drake}}

Our results reframe the Drake Equation's longevity term in a way that addresses the Fermi Paradox with greater nuance. Rather than requiring that intelligent life be exceedingly rare, the Great Silence may simply indicate that technosignatures are intermittent and timing-dependent \citep{Balbi2018-fp}. Aligned with this, recent work suggests that civilizations might commonly face phase transitions or limits to growth that prevent long, unbroken periods of signal emission \citep{Wong2022-hv}. In fact, our bimodal outcomes, with some scenarios yielding repeated collapse and others achieving long-term stability, closely echo this hypothesis. The overall idea is that planetary civilizations may follow one of two paths: either ``asymptotic burnout'', collapsing under the weight of exponential growth, or a ``homeostatic awakening'' in which they deliberately curb expansion to achieve equilibrium. In the latter case, a civilization might survive indefinitely but become hard to detect due to a lack of aggressive expansion or high-power signaling. Our high duty-cycle scenarios like Golden Age (S3) or Out of Eden (S10) resemble this stable, sustainable type, whereas the collapse-prone scenarios like Big Brother (S1) or Sword of Damocles (S6) are akin to burnout cases. The implication is that the rarity of continuous technosignature emitters (high \emph{L} or \emph{D\textsubscript{c}}) in our simulations could reflect a Galactic trend, where long-lived, broadcasting civilizations might be the exception, not the rule.

\textcolor{red}{To formalize this connection, we recall the Drake Equation \citep{Drake1965-wp}, which estimates the number of detectable civilizations in our galaxy as}
\textcolor{red}{\begin{equation}
N = R_*\, f_p\, n_e\, f_l\, f_i\, f_c\, L .
\label{eq:drake}
\end{equation}}
\textcolor{red}{Each parameter represents a step in the emergence of detectable life: \(R_{*}\) is the rate of star formation, \(f_{p}\) the fraction of those stars with planets, \(n_{e}\) the number of habitable planets per system, \(f_{l}\) the fraction of planets where life arises, \(f_{i}\) the fraction where intelligence evolves, \(f_{c}\) the fraction that develop detectable technology, and \(L\) the average duration such civilizations remain detectable. While early terms are increasingly constrained by astrophysical observation, the final term \emph{L} carries outsized importance and remains deeply uncertain. If civilizations exist but are short-lived or frequently silent, \(N\) could remain small despite favorable conditions for life elsewhere.}

An important extension of this idea is that technology can outlast its creators. Our study has focused on the active phase of civilizations (the ON periods), but one must consider that even a civilization which collapses (goes OFF) might leave behind artifacts or long-lived technosignatures. For example, it has been argued that the maximum longevity of a technosignature could be far greater than the civilization's biological lifespan, because machines, probes, or infrastructure can continue operating (or at least remain observable) for millennia or longer \citep{Wright2022-wh}. In this view, Earth and humanity may be poor guides to the true value of \emph{L}, since advanced extraterrestrial technology might persist well beyond the civilization's demise. This is a crucial point for SETI as even extinct or dormant civilizations could be detectable if they've built something durable \citep{Cirkovic2019-oq}. Dyson's classic idea of searching for waste heat from starlight-absorbing megastructures is a case in point. A ``Dyson Sphere'' could remain in orbit around a star for eons, glaring in the infrared, even if the society that built it is long gone \citep{Dyson1960-jf}. Such relic technosignatures (sometimes called ``Dysonian'' technosignatures) broaden our notion of detectability. In the same vein, a dense belt of artificial satellites around a planet (a ``Clarke exobelt'') might be detectable via its subtle transit signature \citep{Socas-Navarro2018-kg}. Other examples include planetary atmospheric anomalies (e.g., chlorofluorocarbon pollution or unusual isotopic ratios in a planet's spectrum) that could signal past industrial activity \citep{Wright2022-wh}. In short, not all technosignatures are ``live broadcasts''. Some might be fossils or afterglows of past epochs of activity.

Our simulations suggest that the traditional interpretation of the Drake Equation's longevity term,\, \emph{L}, as a continuous span of technosignature emission, may be unrealistic. Across a wide range of plausible futures, we find that civilizational activity is frequently interrupted by collapse, resource depletion, or other disruptions. In most scenarios, the emission of technosignatures is not constant but episodic, raising the need for a revised formulation. We propose an effective detectability duration, $L_{\mathrm{eff}} = D_c \times T_{\mathrm{span}}$, where $D_c$ is the mean duty cycle and $T_{\mathrm{span}}$ is the time horizon. This redefinition captures the operational reality of civilizations that switch between active and dormant states. This intermittency reframes how \emph{N}, the number of detectable civilizations, should be interpreted. If $L$ is better understood as $L_{\mathrm{eff}}$, then models that assume continuous detectability may significantly overestimate \emph{N}. Conversely, incorporating duty cycle effects helps align theoretical predictions with empirical null results. A galaxy populated by civilizations with low duty cycles would remain largely silent, even if intelligent life is common.

This framework has two important implications. First, it helps explain the Great Silence. Civilizations may exist in abundance, but if they are detectable only intermittently, the likelihood that their technosignatures overlap with our current observational window is low. The apparent absence of signals may not reflect the rarity of extraterrestrial technology, but rather the rarity of sustained emission. Second, it suggests that long-lived, continuously emitting civilizations are outliers. Most of our simulated futures spent significant portions of time in silent or collapsed states, even if they eventually recovered. If such patterns are typical across the galaxy, then detection becomes a matter of timing: for two intermittent civilizations, the probability of mutual visibility is the product of their duty cycles, which is potentially very small. This stands in contrast to the assumptions embedded in the Kardashev scale \citep{Kardashev1964-gv}, which classifies civilizations by their continuous energy use and implies long-term, uninterrupted detectability. Our simulations instead point to dynamic energy throughput, with civilizations temporarily reaching higher levels, such as approaching Type I, only to regress during periods of collapse or dormancy. These fluctuations reduce their cumulative detectability and challenge the static expectations of the Kardashev framework. Moreover, while our models focus on collapse-driven intermittency, it's worth noting that some civilizations might limit their energy use deliberately (in pursuit of ecological balance or long-term stability, for example), remaining below Type I not due to collapse, but by design.

Our results also help reconcile seemingly contradictory perspectives on civilizational longevity. One view posits that technological civilizations are fundamentally fragile and prone to self-destruction soon after emergence, driven by ecological, behavioral, or structural instabilities \citep{Vinn2024-pi}. Another suggests that once civilizations surpass key thresholds of resilience and coordination, they may persist indefinitely, steadily accumulating across cosmic time \citep{Grinspoon2009-cg}. Both patterns emerge naturally in our simulations. Scarcity-driven, hierarchically governed scenarios frequently enter collapse-recovery cycles, resulting in low duty cycles and limited detectability. In contrast, post-scarcity, distributed scenarios maintain stable resource use and continuous activity, supporting high-duty-cycle technospheres. These differences arise not from exogenous shocks, but from internal sociotechnical structure.

\textcolor{red}{It is worth noting, however, that the Drake Equation's terms beyond $f_l$ carry increasingly problematic hidden assumptions that our duty-cycle framework helps to expose \citep{Cirkovic2004-em}. For the product $f_l \times f_i$ to be pragmatically sound, $f_i$ must represent the fraction of life-bearing planets that develop intelligent life \emph{at least once}, not merely once, acknowledging that intelligence may arise multiple times on a single world. Consequently, $f_c$ must be understood as the fraction of planets with at least one instance of intelligent life on which at least one instance of traceable technology emerges. Most critically, $L$ must refer to the duration of traceability, decoupled from the longevity of the technology's producers and its originating technology. That is, $L$ must also account for the temporal dynamics of remnants left by extinct intelligent builders on or around worlds where new kinds of intelligent builders may not even be producing traceable technologies---a situation of spatial co-existence without temporal co-existence of technology-capable beings. In principle, multiple kinds of intelligent builders could produce different technologies with different degrees of traceability, or technologies that are even antagonistic to one another, although this complicates the picture considerably. Extending the duty-cycle metaphor, the ON-OFF behavior of a technosphere may not be ``gated'' solely by triggers of collapse as a rectangular pulse. Instead, it may represent the reaching and exceeding of a particular detectability threshold---contingent on the observing civilization---through the intentional or unintentional ``constructive'' or ``destructive'' ``interference'' of particular technologies deployed by a single civilization of builders or by multiple spatially co-existent civilizations that may or may not co-exist temporally. For the particular case of the ten scenarios modeled in this paper, we have assumed that humans are the only technology-capable species on Earth, so the issue of multiple kinds of builders does not arise; only formerly collapsed incarnations of human civilization are relevant.}

\textcolor{red}{Recent work by \citet{Balbi2024-bg} provides further theoretical support for the importance of the duty-cycle framework. Applying Lindy's law to technosignature longevities, they show that for energy-intensive technoemissions the longevity distribution follows a power law $\rho_L(L) \propto L^{-\alpha}$ with $\alpha > 1$, strongly disfavoring long-lived technosignatures. Under realistic assumptions ($\alpha \geq 2$), the first detected technosignature is predicted to be relatively short-lived ($L < 10^2$ years at 95\% confidence), rather than the $10^9$-year signals sometimes assumed in optimistic SETI scenarios. This result strengthens the case for our framework: if most detectable technosignatures are intrinsically short-lived, then the combination of Lindy-distributed longevities and low duty cycles (as observed in our collapse-prone scenarios) implies a doubly diminished probability of detection, offering a complementary explanation for the Great Silence.}

\textcolor{red}{\subsection{Duty Cycles of Technological Remnants}\label{duty-cycles-remnants}}

\textcolor{red}{An additional consideration that extends the duty-cycle framework concerns the autonomous behavior of technological remnants. When a civilization collapses, not all of its technological infrastructure necessarily ceases to function. Some systems---particularly those designed for autonomous operation---may continue to operate on their own duty cycles, independent of the civilization that created them. For example, a planetary defense system designed to deflect asteroids might be activated and deactivated at predetermined intervals or in response to triggering events, persisting long after the civilization that built it has collapsed and while a recovering population has no memory of its existence. Similarly, automated beacons, environmental monitoring stations, or orbital maintenance systems could continue cycling through active and dormant states according to their original programming or in response to environmental triggers.}

\textcolor{red}{This creates a layered duty-cycle structure: the remnant technology's ON-OFF cycle is superimposed on the civilization's overall collapse-recovery cycle. During periods when the civilization is OFF (collapsed or recovering), remnant technologies in their ON phase could produce detectable technosignatures, effectively extending the civilization's observable presence beyond its active periods. Conversely, when the civilization recovers and re-enters its ON phase, it may encounter and reactivate dormant remnant systems, creating a form of constructive interference between the two duty cycles. This nuance modifies the persistence term $\tau$ in our extended formulation ($L_{\mathrm{eff}}^{(d)} = D_c \, T_{\mathrm{span}} + N_{\mathrm{ON}} \, \tau$): rather than representing a simple passive decay time, $\tau$ may itself contain active intervals during which remnant technologies produce detectable signatures. The effective detectability window thus becomes the union of the civilization's active periods and the remnant technology's active periods, potentially increasing the overall duty cycle beyond what civilizational activity alone would suggest.}

\textcolor{red}{\subsection{Limitations and Modeling Considerations}\label{limitations}}

Finally, several limitations warrant discussion. Our modeling framework is grounded in Earth-originating scenario typologies, drawing from narrative futures methodologies informed by contemporary geopolitical, ecological, and technological trends. These projections embed assumptions about governance structures, resource constraints, and recovery dynamics that may not extend to non-terrestrial civilizations or to those that do not follow Earth-like developmental pathways. While this grounding enables empirical plausibility, it limits extrapolation beyond our biosphere's trajectory \citep{Rubin2001-sk,Denning2011-fj}. Moreover, our simulations assume bounded planetary systems and do not account for galactic colonization, interstellar migration, or modes of expansion that could alter the spatial and temporal visibility of civilizations \citep{Walters1980-ev,Papagiannis1983-xv}. We also exclude the effects of cultural convergence, information saturation, or technological transcendence (e.g., the emergence of post-biological intelligence), all of which could suppress or qualitatively change technosignature output. Our operational definition of detectability (based on internal sociotechnical continuity) does not differentiate among specific technosignature classes (e.g., radio emissions, waste heat, megastructures), nor does it address detection thresholds or instrumental sensitivity.

\textcolor{red}{Additionally, our assumption that technological capacity increases linearly each year according to a fixed growth rate $r$ is a significant simplification. The many discontinuities implied by the scenarios---not only collapse points but also radical transformation points such as the emergence of artificial general intelligence or breakthroughs in energy harvesting---suggest that a linear growth model may not adequately capture the dynamics of technological change across all regimes. Moreover, the available resource stock $R_0$ is treated as an exogenous parameter, but in practice it is contingent on the technological capacity to identify and utilize resources: a civilization that develops asteroid mining or nuclear fusion effectively increases its accessible resource base, fundamentally altering the growth-depletion dynamics. Recent work by \citet{Haqq-Misra2025-lum}, which shares the same ten-scenario framework employed here, reframes the Kardashev scale as a luminosity limit constrained by thermodynamic efficiency (exergy) and explores exploration-exploitation trade-offs for long-lived technospheres. Their framework highlights that growth trajectories depend on whether a civilization pursues spatial expansion (exploration) or intensive resource use within a fixed domain (exploitation), a distinction that our linear growth model does not capture. Future work could incorporate exergy constraints and non-linear growth functions informed by their luminosity-limit model, as well as feedback between technological capacity and accessible resource stocks, to produce more physically grounded projections of civilizational trajectories.}

In summary, our results can partially explain the Great Silence: if detectability is often sporadic, then long periods without contact may be the expected baseline. Intermittency may be typical and recognizing that could refine our expectations and expand the scope of SETI.
